\documentclass[11pt]{article}

% ---------------------------------------------------------------------------
%  L1 — A diagonal law for row-local lattices, and the mod-3 arithmetic of the
%  king instance.
%  Category L (entirely machine-written); see paper/README.md and
%  docs/publication-split.md.
%
%  MERGED 2026-08-23.  L2-ternary-spine.tex was folded in as Part II and
%  withdrawn as a manuscript.  This paper's abstract already advertised the
%  spine cubic W^3 = W^2 + t as its own third consequence, and every theorem in
%  L2 was conditional on Theorems A and C here -- L2 cited this paper eight
%  times and spent its entire Setup restating the two results it imported.
%  docs/l-corpus-contraction.md is the record.
%
%  NOVELTY: no collision on either half.  Part I's N1-N6 sweeps stand; Part II
%  was swept 2026-08-18 (docs/priority-passes-2026-08-18.md) with no collision
%  on the spine cubic, the digit product or the Smith normal form count.
%
%  Source material, all in this repository:
%    docs/proofs/universal-diagonal-law.md   Theorems A and B, the instances
%    docs/proofs/diagonal-law.md             the king case, written out first
%    docs/proofs/grand-form.md               Theorem C and its corollaries
%    docs/proofs/T-n-nm1.md                  P_1(n) = 25n - 45 from scratch
%    docs/proofs/T-n-nm2-and-general.md      the king k = 2 hand derivation
%    results/hex-diagonal-law.md             the b = 2 instance
%    results/polyiamond-diagonal-law.md      the periodic extension
%    results/diagonal-closed-forms.md        the literature check on T(n,n-k)
%    results/defect-gas.md                   the king cluster weights
%    results/ternary-spine.md                Part II, the mod-3 arithmetic
% ---------------------------------------------------------------------------

% ---------------------------------------------------------------------------
%  paper/shared/preamble.tex — packages and macros common to every manuscript
%  in this directory.  Factored out of paper/polyplets-report.tex and
%  paper/technical-report.tex, which between them had two copies of the same
%  twenty lines.
%
%  technical-report.tex does NOT \input this file and is not to be edited to do
%  so: it is jasonp's prose and is read-only to the machine (paper/README.md,
%  "Who may edit what").  The duplication there is deliberate.
%
%  Assumes \documentclass[11pt]{article} has already been issued.
% ---------------------------------------------------------------------------
\usepackage[margin=1in]{geometry}
\usepackage{amsmath,amssymb}
\usepackage{amsthm}
\usepackage{booktabs}
\usepackage{graphicx}
\usepackage{numprint}
\usepackage{microtype}
\usepackage[table]{xcolor}
\usepackage[hidelinks]{hyperref}
\usepackage{flafter}   % a float never appears before the text that introduces it
\usepackage{float}     % [H] pins tables/figures exactly where written
\usepackage[font=small,labelfont=bf,labelsep=period]{caption}
\setlength{\intextsep}{16pt plus 3pt minus 3pt}

\npthousandsep{,}
\npfourdigitsep
\newcommand{\num}[1]{\numprint{#1}}
\newcommand{\g}{\cellcolor{gray!15}}

% Theorem environments.  One counter shared across kinds, so that a reference
% like "Theorem 4" is unambiguous in a paper that also has lemmas and
% propositions.
\theoremstyle{plain}
\newtheorem{theorem}{Theorem}
\newtheorem{proposition}[theorem]{Proposition}
\newtheorem{lemma}[theorem]{Lemma}
\newtheorem{corollary}[theorem]{Corollary}
\theoremstyle{definition}
\newtheorem{definition}[theorem]{Definition}
\newtheorem{example}[theorem]{Example}
\theoremstyle{remark}
\newtheorem{remark}[theorem]{Remark}
\newtheorem{conjecture}[theorem]{Conjecture}
\newtheorem{openproblem}[theorem]{Open Problem}

% Repo-path citations.  Every one of these must resolve in the tree that ships
% with the paper; `make gate-citations` checks the markdown, and
% paper/verify_claims.py checks the manuscripts.
\newcommand{\repo}[1]{\texttt{#1}}

% OEIS links.
\newcommand{\oeis}[1]{\href{https://oeis.org/#1}{#1}}

% Growth constants, so a rename is one edit.
\newcommand{\lam}{\lambda}
\newcommand{\muH}{\mu_H}

% ---------------------------------------------------------------------------
%  paper/shared/disclosure.tex — the two authorship disclosure blocks.
%
%  docs/publication-split.md §1 fixes the rule these implement:
%
%    "Under no circumstances will there be any ambiguity about whether a paper
%     is wholly my work (none), merely my prose (some) or entirely LLM-authored
%     (some)."   — jasonp, 2026-08-07
%
%  There are exactly two blocks and they live in exactly one file, so that a
%  paper cannot quietly soften its own disclosure.  The fixed sentences take no
%  arguments.  The ONE thing a paper supplies is the per-result verification
%  ledger, because that is the part that differs honestly between papers — and
%  §1 requires it to be per result, not in aggregate.
%
%  Do not add a \Pdisclosure or \Ldisclosure variant with an optional softening
%  argument.  If a paper does not fit one of these two blocks, the paper is
%  wrong, not the block.
% ---------------------------------------------------------------------------

\newsavebox{\disclosurebox}

\newenvironment{disclosureframe}
  {\begin{center}\begin{lrbox}{\disclosurebox}\begin{minipage}{0.93\textwidth}\small}
  {\end{minipage}\end{lrbox}\fbox{\usebox{\disclosurebox}}\end{center}}

% --- P: jasonp's prose -------------------------------------------------------
% Byline: Jason Parker, alone.  Argument: the per-result ledger of what the
% machine did.
\newcommand{\Pdisclosure}[1]{%
\begin{disclosureframe}
\textbf{Authorship.}\enspace Every sentence of this paper was written by its
named author. He has read every proof it states and believes each one, and he
is responsible for the correctness of what it claims.

\smallskip
\textbf{Machine contribution.}\enspace #1
\end{disclosureframe}}

% --- L: entirely machine-written --------------------------------------------
% Argument: the per-result verification ledger — for each result, what warrants
% it (Lean kernel, exact certificate, or labelled measurement) and how much of
% it a human has checked.
\newcommand{\Ldisclosure}[1]{%
\begin{disclosureframe}
\textbf{Authorship.}\enspace This document was written end to end by a large
language model. That includes its mathematics, its prose, its tables and its
\LaTeX{}. No sentence of it was written by a human. The mathematics it reports
was likewise derived by machine.

\smallskip
\textbf{Human role.}\enspace The person who directed this work chose what to
compute, chose what to publish, and is responsible for the decision to release
this document. Except where the ledger below says otherwise, that person has
not personally checked the arguments here, and this disclosure is not to be
read as an endorsement of them.

\smallskip
\textbf{What warrants each result.}\enspace #1
\end{disclosureframe}}

% --- Draft banner ------------------------------------------------------------
% For an L-paper whose verification ledger is not yet settled.  Loud on purpose:
% an L-paper without a truthful ledger must not be mistaken for a finished one.
%
% \firstdraftbanner is the standard wording, which five papers previously
% carried character-for-character in their own preambles.  It lives here for the
% same reason the disclosure blocks do: identical text in seven places drifts,
% and a paper must not be able to soften it by editing its own copy.  The
% optional argument is for a gate specific to one paper, appended verbatim; a
% paper needing different wording (L6, compute-gated) calls \draftbanner direct.
\newcommand{\firstdraftbanner}[1][]{%
\draftbanner{This is a first draft. The verification ledger below records what
warrants each result \emph{mechanically}; the column that records what a human
has read is empty, deliberately, because nobody has read it yet. Do not
circulate until that column says something true.#1}}

\newcommand{\draftbanner}[1]{%
\begin{center}
\fcolorbox{red!60!black}{yellow!15}{%
  \begin{minipage}{0.93\textwidth}\small
  \textbf{DRAFT — NOT FOR CIRCULATION.}\enspace #1
  \end{minipage}}
\end{center}}


\newcommand{\Wp}{W_{\mathrm{pair}}}
\newcommand{\Z}{\mathbb{Z}}
\newcommand{\Q}{\mathbb{Q}}
\newcommand{\F}{\mathbb{F}}
\newcommand{\vthree}{v_3}

\title{\textbf{A diagonal law for row-local lattices}\\[4pt]
\large Polynomial closed forms for lattice-animal counts at fixed
height defect, and the mod-$3$ arithmetic of the king instance}
\author{[byline: jasonp's call --- \repo{docs/publication-split.md} \S1]}
\date{August 2026}

\begin{document}
\maketitle

\firstdraftbanner

\Ldisclosure{%
\textbf{Part I --- the law.}
\emph{Theorem~\ref{thm:A} (shape, onset, integrality)} --- proved below. The
king instance ($b=3$) is formalized in Lean~4 against mathlib
(\repo{polyplets/Polyplets/Universal/}: \texttt{universal\_shape\_d},
\texttt{universal\_shape}, \texttt{universal\_shape\_production},
\texttt{universal\_production\_int\_all}) and depends on standard axioms only,
with no \texttt{sorry} anywhere in the development. Machine checks of the
general statement: \repo{experiments/universal\_law\_check.py} (square),
\repo{experiments/diagonal\_law\_proof\_check.py} (king),
\repo{experiments/hex\_gas.py} (hex).
\emph{Theorem~\ref{thm:B} (the spine)} --- paper proof only, not formalized;
both sharpness witnesses are machine computations
(\repo{experiments/universal\_pair\_weights.py}).
\emph{Theorem~\ref{thm:C} (the product form)} --- paper proof below; the king
instance is Lean, standard axioms only (\texttt{grand\_form},
\texttt{grand\_form\_prod}), and the per-level production pinning for
$k \le 18$ carries one named \texttt{native\_decide} leaf per weight card.
\emph{The closed forms of Table~\ref{tab:machine}} --- each is a consequence of
Theorem~\ref{thm:A} applied to $k+1$ enumerated values, and each survives 5--11
withheld values that were not used to fit it.
\emph{Corollary~\ref{cor:a308359}} --- deduced from Theorem~\ref{thm:A}; the
enumeration behind it independently reproduces the case Mathar records as
proved.
\textbf{Part II --- the mod-$3$ arithmetic.}
\emph{Theorem~\ref{thm:snf-3power} (all invariant factors are $3$-powers)} ---
unconditional, two lines, proved below.
\emph{Propositions~\ref{prop:polya} and \ref{prop:GH}} --- unconditional given
the diagonal law.
\emph{Theorems~\ref{thm:digit}, \ref{thm:oddspine}, \ref{thm:evenspine},
\ref{thm:activation}, \ref{thm:snf}} --- conditional on the diagonal law
(Theorem~\ref{thm:A}, proved in Part I) together with the three congruences of
\S\ref{sec:ladder}, each derived symbolically from the cluster master equation. All
five are additionally verified against an independently enumerated triangle:
$342$ in-regime cells for the digit product, all $k \le 17$ for the spine
theorems, and every $N \le 36$ for the Smith normal form.
\emph{Theorem~\ref{thm:deficit2} (the deficit-$2$ residue)} --- proved by
Lagrange--B\"urmann over $\Z/27$ followed by an exact polynomial division on
the mod-$27$ master curve; the division is a finite exact computation.
\emph{Theorem~\ref{thm:d3} (the deficit-$3$ residue)} --- proved on the frame
stated with it: (A1) and (A4) are proved and independently cross-checked, while
(A2) and (A3) are symbolic derivations verified against computed data to $k \le 17$
rather than proofs. Everything above the frame is exact finite algebra.
\emph{\S\ref{sec:sleeve} (sleeve zeros)} --- labeled measurement, complete for
$n \le 40$, with no closed form claimed \emph{except} at deficit $3$, where
Theorem~\ref{thm:d3} predicts them for all $k$.
\emph{Novelty} --- pass run 2026-08-18, no collision found on the spine cubic,
the digit product or the Smith normal form count; the two bodies of machinery a
reader will think of are named in \S\ref{sec:related}, with why neither applies.
\emph{Human verification:} none, as of this draft, on either part.}

\begin{abstract}
\noindent
\textbf{Part I.}
Let $T(n,H)$ count the fixed lattice animals of $n$ cells whose bounding box has
height exactly $H$. We prove that for the \emph{row-local} lattices, those whose
adjacency moves the row index by at most one,
the diagonals of this triangle have closed forms: for every $k \ge 0$ there is a
polynomial $q_k$ of degree exactly $k$ with
\[
  T(H+k,\,H) \;=\; q_k(H)\cdot b^{H}
  \qquad\text{for all } H \ge k+1,
\]
where $b$ is the number of ways the lattice steps from one row to the next. The
onset $H \ge k+1$ is proved, not observed. The leading coefficient of $q_k$ is
$\Wp^{\,k}/k!$, where $\Wp$ is a single enumerable constant of the lattice.
Polyominoes ($b=1$), polyhexes ($b=2$) and polyplets ($b=3$) are three
instances of one theorem.

Three consequences. First, $k+1$ enumerated values of a diagonal now
\emph{determine} it, which turns a fitted formula into a proved one; we tabulate
the resulting closed forms. Second, the
conjecture recorded at OEIS \oeis{A308359} follows: the width-defect-$2$
diagonal of the fixed polyomino triangle is $8n^2-51n+86$ for $n \ge 5$, degree
and onset included. Third, modulo a prime $p$ dividing $b$ the
generating function of the law satisfies one cubic, $W^3 = W^2 + t$, that does
not depend on the lattice: the lattice chooses only the prime and a scaling.

We also record where the theorem stops. Polyiamonds fail the hypothesis, their
adjacency being orientation-dependent, so that no single step set exists; the
conclusion nevertheless survives in periodic form, which we report as
measurement rather than theorem.

\textbf{Part II} takes the third of those consequences seriously on one
lattice. Let $T_N = [T(i,j)]_{i,j \le N}$ be the leading block of the king
triangle. It is lower triangular with determinant $3^{N(N-1)/2}$, so its Smith
normal form over $\Z$ is $\mathrm{diag}(3^{e_1},\dots,3^{e_N})$; the number of
nontrivial invariant factors is $\lceil (N-1)/3\rceil$, which does not follow
from the determinant. We derive it from the arithmetic of the triangle modulo
$3$, using Part I's law and product form together with three congruences read
off the cluster master equation.

Writing $n = \sum_i n_i 3^i$ in base $3$, the diagonal polynomials satisfy
\[
  P_k(n) \;\equiv\; [y^k]\prod_i W\!\left(y^{3^i}\right)^{n_i} \pmod 3 ,
\]
with $W$ the spine of Theorem~\ref{thm:B}, so the triangle modulo $3$ is
$3$-automatic with an explicit automaton. The cubic admits the parametrisation
$t = x/(1-x)^3$, $W = 1/(1-x)$, in which the Lagrange--B\"urmann correction term
vanishes identically in characteristic $3$, so the two boundary cases close:
column $H$ vanishes modulo $3$ in every row below $\lfloor 3H/2\rfloor$ and is a
unit in that row. Since $H \mapsto \lfloor 3H/2\rfloor$ is strictly increasing,
the surviving columns stand in echelon position, and counting the vanishing ones
gives the invariant-factor count.

Read along rows, the last nonzero residue of row $n$ is $1$, $1$, $2$
according to $n \bmod 3$, the third case being the deficit-$2$ cell, which we
treat separately by Lagrange--B\"urmann over $\Z/27$. We give a complete census
for $n \le 40$ of the cells whose $3$-adic valuation exceeds the deficit the law
forces on them; a closed formula for those exceptions is open. The cubic lifts:
$H^3 - H^2 \equiv 25t \pmod{27}$ with $25$ the single-defect weight of the
cluster decomposition, and $W$ returns as the level-$3$ correction.
\end{abstract}

\tableofcontents

\part*{Part I --- the law}
\addcontentsline{toc}{part}{Part I --- the law}

\section{Introduction}

The number of fixed lattice animals with $n$ cells has no known closed form on
any lattice of interest, and there is good reason to think none exists: the
anisotropic generating function of the polyplet height triangle is provably
not D-finite, and the same argument applies to polyominoes and polyhexes
(companion paper L4, \emph{An arithmetic obstruction to D-finiteness for
lattice-animal height generating functions}).
Stratifying by height does not repair that; it exposes structure the row sums
hide.

Write $T(n,H)$ for the number of fixed animals with $n$ cells whose bounding box
has height exactly $H$. The rows of this triangle are the objects one normally
wants; the \emph{diagonals}, got by fixing the \emph{defect} $k = n - H$ and
letting $H$ grow, are the objects that turn out to be tractable. Empirically, on every
lattice anyone has tabulated, they are polynomial times exponential. This paper
proves that, in the following form.

\begin{quote}
For every row-local lattice and every $k \ge 0$ there is a polynomial $q_k$ of
degree $k$ with $T(H+k,H) = q_k(H)\,b^H$ for all $H \ge k+1$, where $b$ counts
the ways of stepping from a row to the next.
\end{quote}

\subsection{Relation to existing work}
\label{sec:priorart}

Two searches came back empty: one for a height-diagonal law for lattice-animal
triangles, one for published closed forms for $T(n,n-k)$. Both are recorded in
\repo{results/diagonal-closed-forms.md} and \repo{results/novelty-sortie.md}.
We claim the theorem, not priority, and cite the nearest published relative.

That relative measures the defect in \emph{dimension} rather than in height, on
polycubes. Barequet, Barequet and Rote~\cite{barequetBarequetRote2010} show that
for fixed $n$ the number of polycubes is a polynomial in the dimension $d$;
Barequet and Shalah~\cite{barequetShalah2017} give, for general $k$, the shape
\[
  2^{\,n-2k+1}\, n^{\,n-2k-1}\,(n-k)\, h_k(n)
\]
for the count of $n$-cell polycubes proper in $n-k$ dimensions, with $h_k$
polynomial. That is the same statement shape as ours (fix a defect parameter, get
polynomial times exponential, with the degree governed by the defect) for
one lattice family and one defect. Theorem~\ref{thm:A} quantifies over every
row-local lattice instead, and is proved once for all of them.

The one place where our theorem settles something on the record is
Corollary~\ref{cor:a308359}, which proves a conjecture standing in OEIS
\oeis{A308359}~\cite{a308359} since 2019.

\section{Row-local lattices}
\label{sec:class}

\begin{definition}[row-local lattice]
\label{def:rowlocal}
A \emph{row-local lattice} has cells indexed by $\Z^2$ and a
translation-invariant symmetric adjacency relation $A \subseteq \Z^2$ (a set of
offsets, with $u \sim v$ iff $v - u \in A$) satisfying
\begin{description}
\item[(L)] $A \subseteq \{(dx,dy) : |dy| \le 1\}$ --- adjacency changes the row
  index by at most one;
\item[(R)] the within-row offsets $\{dx : (dx,0) \in A\}$ are finite and
  contain $dx = \pm 1$ --- a run of consecutive cells in one row is connected;
\item[(U)] the \emph{up-offset set} $D := \{dx : (dx,1) \in A\}$ is finite and
  nonempty. Write $b := |D|$, the \emph{drift count}.
\end{description}
An \emph{animal} is a finite nonempty connected set of cells, taken up to
translation; $T(n,H)$ counts the animals with $n$ cells occupying exactly $H$
consecutive rows.
\end{definition}

Condition (L) is the whole of the geometric content --- it makes a row with a
single cell a cut, which is Lemma~\ref{lem:sep}. Conditions (R) and (U) are finiteness bookkeeping.

\begin{example}[the three instances]
\label{ex:instances}
\leavevmode
\begin{itemize}
\item \textbf{Square lattice.} $A = \{(\pm1,0),(0,\pm1)\}$, so $D = \{0\}$ and
  $b=1$. Animals are fixed polyominoes, \oeis{A001168}.
\item \textbf{Polyhexes.} $D = \{-1,0\}$, $b=2$. Cells are hexagons, each with
  six neighbors; $T$'s row sums are $1, 3, 11, 44, 186, 814, 3652, 16689,
  \dots$ = \oeis{A001207}, fixed hexagonal polyominoes.
\item \textbf{Polyplets (the king lattice).} $A$ is all eight compass offsets, so
  $D = \{-1,0,1\}$ and $b=3$. Row sums are \oeis{A006770}.
\end{itemize}
\end{example}

\begin{remark}[a name collision]
\label{rem:names}
The centres of the $b=2$ cells form the triangular \emph{point} lattice, so the
literature sometimes calls this the triangular lattice --- Madras's
``Tri''~\cite{madras1999}, for instance. That is the six-neighbor object and it
gives polyhexes. \emph{Tiling by triangles} is a different object: each
equilateral cell has three edge-neighbors, the counts are
$2, 3, 6, 14, 36, 94, 250, 675, 1838, \dots$ = \oeis{A001420}, and OEIS files
these polyiamonds under ``the 2-dimensional hexagonal lattice''.
Definition~\ref{def:rowlocal} covers the first and not the second; see
Section~\ref{sec:polyiamond}.
\end{remark}

\subsection{The decomposition}

Fix a row-local lattice and read an animal of height $H$ row by row; every one of
its $H$ rows is nonempty. Call a row with exactly one cell a \emph{walk row}, and
call a maximal run of rows with two or more cells a \emph{cluster}. A cluster's
\emph{type} is its stack of row sizes $(s_1,\dots,s_\ell)$ with every
$s_i \ge 2$; its \emph{surplus} is $k_c = \sum_i (s_i - 1)$ and its \emph{length}
is $\ell_c$. Since every cluster row contributes at least one to the surplus,
\begin{equation}
  \ell_c \;\le\; k_c .
  \label{eq:elllek}
\end{equation}
Inequality~\eqref{eq:elllek} is used at three points below: the degree bound of
Step~4, the onset $H \ge k+1$, and the collapse of the mod-$p$ series in
Section~\ref{sec:spine}.

Four families of weights count the configurations of a cluster relative to its
contacts:
\begin{itemize}
\item the \emph{interior} weight $W_c$: configurations of a cluster of type $c$
  together with the single walk cell below it (held fixed) and the single walk
  cell above it (free), the whole being connected;
\item the \emph{edge} weights $W^b_c$ and $W^t_c$: the same with the cell below
  (respectively above) absent, for clusters that touch the bottom (respectively
  top) of the animal. $W^t$ of a type is $W^b$ of its reversal;
\item the \emph{pure} weights $W^p_c$: animals consisting of a single cluster and
  nothing else.
\end{itemize}

All four weight families are finite, and each surplus admits only finitely
many cluster types: by (R) and (U) a connected set of cells has bounded
$x$-spread, so clusters of bounded surplus range over finitely many
configurations once $x$-translation is quotiented out. The sums in
Proposition~\ref{prop:chain} are therefore formal power series.

\section{The diagonal law}
\label{sec:A}

\begin{theorem}[diagonal law]
\label{thm:A}
Let $L$ be a row-local lattice with drift count $b$. For every $k \ge 0$ there is
a polynomial $q_k \in \Q[H]$ of degree at most $k$ such that
\[
  T(H+k,\,H) \;=\; q_k(H)\cdot b^{H}
  \qquad\text{for every } H \ge k+1 .
\]
Equivalently, setting $P_k(n) := b^{\,1+2k}\, q_k(n-k)$,
\[
  T(n,\,n-k) \;=\; P_k(n)\cdot b^{\,n-1-3k}
  \qquad\text{for every } n \ge 2k+1 ,
\]
and $P_k$ takes an integer value at every integer.
\end{theorem}

The proof is five steps. None of them uses the value of $b$.

\subsection*{Step 1: a walk row is a cut}

\begin{lemma}[separation]
\label{lem:sep}
A set of cells occupying rows $r_1 < \dots < r_H$, every row nonempty, is
connected if and only if each segment between consecutive walk rows (endpoints
included) is connected.
\end{lemma}

\begin{proof}
By (L), adjacency changes the row index by at most one, so any path from a cell
strictly below row $r$ to a cell strictly above row $r$ must contain a cell of
row $r$. If $r$ is a walk row, that cell is unique.

($\Leftarrow$) Consecutive segments overlap exactly in a walk cell, so
connectedness glues.

($\Rightarrow$) Take two cells inside one segment $[r,r']$ and a path between
them in the animal. Any excursion of the path below $r$ leaves and re-enters
through the unique walk cell of row $r$, and likewise above $r'$; excising each
excursion leaves a path inside the segment.
\end{proof}

\subsection*{Step 2: the chain identity}

Let $y$ mark surplus and $z$ mark rows, and set
$F(y,z) = \sum_{n,H} T(n,H)\, y^{\,n-H} z^{H}$.

\begin{proposition}[chain identity]
\label{prop:chain}
\[
  F \;=\; E_b\,(1-S)^{-1} E_t \;+\; P,
\]
exactly, where
\[
  S = b\,z + \sum_c W_c\, y^{k_c} z^{\ell_c+1},\quad
  E_b = z\Bigl(1 + \sum_c W^b_c\, y^{k_c} z^{\ell_c}\Bigr),\quad
  E_t = 1 + \sum_c W^t_c\, y^{k_c} z^{\ell_c},
\]
and $P = \sum_c W^p_c\, y^{k_c} z^{\ell_c}$, all sums over cluster types $c$.
\end{proposition}

\begin{proof}
By Lemma~\ref{lem:sep} every animal decomposes \emph{uniquely} as either a single
pure cluster (the $P$ term) or as an optional bottom-edge cluster, then an
alternating chain of walk steps and clusters cut at its walk rows, then an
optional top-edge cluster. Configurations of each piece are counted relative to
the walk cell it shares with its neighbor, so the choices multiply; rows and
surplus add, so the monomials multiply. A plain walk step contributes $b z$,
since by (U) there are exactly $b$ ways to place the next single cell above the
current one; a cluster together with the walk row above it contributes
$W_c y^{k_c} z^{\ell_c+1}$. Summing the geometric series in $S$ gives the
identity. Uniqueness of the decomposition and invariance under $x$-translation
make it exact rather than an inequality.
\end{proof}

For the king lattice this identity was checked against an independently
enumerated triangle: exact agreement at all 40 cells with $k \le 3$, $H \le 10$,
and per-order agreement to $H = 33$
(\repo{experiments/diagonal\_law\_proof\_check.py}).

\subsection*{Step 3: rational structure}

Write $\sigma := S - bz$, so that every term of $\sigma$ has $y$-order at least
one. Then $(1-S)^{-1} = \sum_{m \ge 0} \sigma^m (1-bz)^{-(m+1)}$, and extracting
$[y^k]$ kills every $m > k$. Hence

\begin{proposition}
\label{prop:rational}
$\displaystyle [y^k]\,F \;=\; \frac{R_k(z)}{(1-bz)^{\,k+1}}$
with $R_k \in \Z[z]$ and $\deg R_k \le 2k+1$.
\end{proposition}

\begin{proof}
Only the degree count needs saying, and it is \eqref{eq:elllek} four times. An
$E_b$ factor at surplus $k_b$ has $z$-degree at most $k_b + 1$; an $E_t$ factor at
surplus $k_t$ has $z$-degree at most $k_t$; a term of $\sigma^m$ at total surplus
$k_m$ has $z$-degree $\sum_i (\ell_i + 1) \le k_m + m$, and carries a factor
$(1-bz)^{\,k-m}$ to reach the common denominator. With
$k_b + k_m + k_t = k$ the total is at most
$(k_b+1) + (k_m+m) + k_t + (k-m) = 2k+1$. The pure part contributes $z$-degree at
most $k$ against a denominator $(1-bz)^{k+1}$, which is within the bound.
\end{proof}

For the king lattice the bound is attained: $\deg R_k$ is exactly $1,3,5,7$ for
$k = 0,1,2,3$.

\subsection*{Step 4: partial fractions and the onset}

Change basis in the numerator: write
$R_k(z) = \sum_{j \ge 0} a_j (1-bz)^{j}$, which is legitimate over $\Q$ and has
$a_j$ equal to the coefficients of $R_k\bigl((1-w)/b\bigr)$, hence
$a_j \in b^{-(2k+1)}\Z$. Splitting the sum at $j = k+1$,
\[
  [y^k]F \;=\; \underbrace{\sum_{j \le k} a_j (1-bz)^{\,j-k-1}}_{\text{poles}}
             \;+\; \underbrace{\sum_{j \ge k+1} a_j (1-bz)^{\,j-k-1}}_{=:\,D(z)} ,
\]
and $D$ is a \emph{polynomial} of degree at most $(2k+1)-(k+1) = k$, by
Proposition~\ref{prop:rational}. Taking $[z^H]$ with $H > \deg D$, that is with
$H \ge k+1$, the polynomial part contributes nothing and
\begin{equation}
  T(H+k,\,H) \;=\; \Bigl(\sum_{i=1}^{k+1} a_{k+1-i}\binom{H+i-1}{\,i-1\,}\Bigr) b^{H}
  \;=:\; q_k(H)\, b^{H},
  \label{eq:qk}
\end{equation}
a polynomial in $H$ of degree at most $k$ times $b^H$, exactly.

The onset is therefore not an artifact of the method: $H \ge k+1$ is precisely
the reach of $D$, and $\deg D \le k$ traces back through
Proposition~\ref{prop:rational} to $\ell_c \le k_c$.

\subsection*{Step 5: integrality}

From \eqref{eq:qk},
\[
  P_k(n) \;=\; b^{\,1+2k} q_k(n-k)
  \;=\; \sum_{i=1}^{k+1} b^{\,1+2k} a_{k+1-i}\binom{n-k+i-1}{\,i-1\,}.
\]
Each binomial is an integer at every integer $n$, and
$v_p(a_j) \ge -(2k+1)v_p(b)$ for every prime $p \mid b$ with no other primes
appearing in a denominator; so every coefficient $b^{1+2k}a_{k+1-i}$ is an
integer and $P_k(n) \in \Z$. This proves Theorem~\ref{thm:A}. \qed

\begin{remark}[integer values, not integer coefficients]
$P_k$ is integer-\emph{valued}; its coefficients in the monomial basis are not
integers for any $k \ge 2$, since by Corollary~\ref{cor:lead} the leading one is
$\Wp^k/k!$ --- for the king lattice, $P_2$ leads with $625/2$. What is true, and
what an implementation should store, is the Newton-basis form implicit in
\eqref{eq:qk}, equivalently the statement that $k!\,P_k \in \Z[n]$. That
statement is a theorem (integer-valuedness plus $\deg \le k$) and is formalized.
$k!$ is not the \emph{minimal} denominator; the minimal one has been worked out
for the king lattice, but it is a fact about a \emph{basis choice} rather than
about the lattice: in the Newton basis every coefficient is an integer by
construction, and there is no denominator left to describe. So it is quoted
nowhere here and nothing below depends on it.
\end{remark}

\section{Sharpness}
\label{sec:sharp}

Theorem~\ref{thm:A} bounds the degree. The bound is attained, and the leading
coefficient is a single enumerable constant of the lattice.

\begin{definition}[the pair weight]
\label{def:pair}
$\Wp$ is the interior weight $W_c$ of the unique cluster type of surplus $1$: a
single row of two cells, counted together with the walk cell fixed below it and
the walk cell free above it.
\end{definition}

Precisely: $\deg q_k = k$ exactly, the coefficient of $n^k$ in $P_k$ is
$\Wp^{\,k}/k!$, and in particular $[n^1]P_1 = \Wp$. That is
Corollary~\ref{cor:lead}, proved in Section~\ref{sec:grand} once the product
form is available. It is written out for the king lattice in full, and the
argument is uniform in the weights, with $25$ replaced by $\Wp$ throughout.

\begin{table}[H]
\centering
\begin{tabular}{l r r l l}
\toprule
lattice & $b$ & $\Wp$ & $P_1(n)$ & onset \\
\midrule
square    & 1 & 4  & $4n-8$   & $n \ge 2k+1$ \\
polyhex   & 2 & 9  & $9n-15$  & $n \ge 2k+1$ \\
king      & 3 & 25 & $25n-45$ & $n \ge 2k+1$ \\
\bottomrule
\end{tabular}
\caption{The three instances at $k=1$.}
\label{tab:instances}
\end{table}

Since $\Wp \ge 4$ on every lattice in the class, the degree bound of
Theorem~\ref{thm:A} is never slack at the top.

\begin{remark}[$4, 9, 25$ are not squares]
\label{rem:notsquares}
The first three values are $2^2$, $3^2$ and $5^2$, which invites the reading that
$\Wp$ is always a perfect square. It is false. Counting with the column gap
uncapped, which is required once $\max|dx| > 1$, gives $\Wp = 58$ and $114$ for
the interval step sets with $b = 4, 5$, and $\Wp = 57$ for $D = \{-2,0,2\}$. The
squares are a $b \le 3$ artifact. What survives is the weaker congruence used in
Section~\ref{sec:spine}: the solid-pair contribution is $(b+1)^2 \equiv 1$
modulo any $p \mid b$.

Interval step sets happen to satisfy $\Wp(b) = b^3 - b(b+1)/2 + 4$, verified for
$b \le 8$ by gadget count and $b \le 5$ end to end. We use that only as a
convenient way to \emph{generate} witnesses for Theorem~\ref{thm:B}'s sharpness,
and we do not present it as a result: exactly three members of the class are
objects anyone studies, and all three are in Table~\ref{tab:instances}.
\end{remark}

\paragraph{Onset sharpness.} Theorem~\ref{thm:A} proves validity from
$n = 2k+1$. That the formula \emph{fails} at $n = 2k$ is a non-cancellation
statement, needing $\deg D = k$ exactly. On the king lattice it holds for every
$k \ge 1$. Companion paper L8 identifies the top coefficient of $R_k$ with the
depth-one below-onset defect $D_1(k) = T(2k,k) - P_k(2k)\,3^{-(k+1)}$ and derives,
by the kernel method and with no series data, an irreducible quartic $\Phi(x,W)$
annihilating $N(x) = \sum_{k \ge 1} 3^{k+1} D_1(k)\,x^k$. The coefficients
$N_k = 3^{k+1}T(2k,k) - P_k(2k)$ are integers, because $P_k$ has degree $k$ and
takes the integer values $T(n,n-k)\,3^{3k+1-n}$ at the $k+1$ consecutive integers
$n = 2k+1, \dots, 3k+1$, hence integer values on all of $\mathbb{Z}$. Reduce
$\Phi$ modulo $3$. It factors:
\[
  \Phi(x,W) \;\equiv\; 2\,(W-1)^3\bigl((1+x)W - x\bigr)
  \qquad\text{in } \F_3[x,W].
\]
Since $\F_3[[x]]$ is an integral domain, the reduction of $N$ kills one factor,
and $W \equiv 1$ would force $N_1 \equiv 0$ against $N_1 = 4$. So the reduction
is $x/(1+x)$, giving $N_k \equiv (-1)^{k+1} \pmod 3$, hence $N_k \ne 0$, hence
$D_1(k) \ne 0$, for every $k \ge 1$. That is $\deg R_k = 2k+1$, which is onset
sharpness. Corollary~\ref{cor:a308359} still uses it only as a control.

The generating function is algebraic and not rational: the all-pairs weight
family is not C-finite, refuted at $\ell = 16$, because the walk between the pair
cells has unbounded gap. The argument rests on L8's identity expressing $D_1$
through the all-pairs cluster families, derived from the chain identity and
machine-verified at $19$ exact cells but not written out as a standalone proof.
The same construction with the drift set as a parameter
(\repo{experiments/depth1\_parametric.py}) gives the square and polyhex
instances: on the square lattice $N(x) = x/(1+x)$ outright, so
$D_1(k) = (-1)^{k+1}$, and the polyhex equation is a quadratic that factors
modulo $2$ as $(W+1)((1+x)W+x)$, so $N_k$ is odd. The onset is sharp on all
three row-local instances.

\section{The product form}
\label{sec:grand}

Theorem~\ref{thm:A} says each diagonal is a polynomial. It says nothing about how
the diagonals relate to each other, and each new $k$ costs exactly two new
rational constants.

\begin{theorem}[product form]
\label{thm:C}
There are $C, \mu \in \Q[[y]]$ with $C(0) = 1/b$ and $\mu(0) = b$ such that
\[
  T(H+k,\,H) \;=\; [y^k]\bigl(C(y)\,\mu(y)^{H}\bigr)
  \qquad\text{for every } k \ge 0 \text{ and } H \ge k+1 .
\]
Equivalently there are rational constants $a_j, b_j$ $(j \ge 1)$, with
$(a_j,b_j)$ determined by the cluster weights of surplus at most $j$, such that
\[
  P_k(n) \;=\; [y^k]\exp\Bigl(\sum_{j \ge 1} (a_j + b_j n)\, y^{j}\Bigr)
  \qquad\text{for every } k \ge 0,\ n \ge 2k+1 .
\]
\end{theorem}

Everything happens in $\Q[[y]]$, $y$-adically, with no analytic or convergence
question anywhere in it, so the proof covers all $k$ at once. We
sketch the four steps; each is elementary and exact per order.

\paragraph{Step 1: Weierstrass preparation by hand.} There is a unique pair
$(z^*, u)$ with $z^* \in 1/b + y\Q[[y]]$, $u = \sum_j u_j(z) y^j$ with
$\deg u_j \le j$ and $u_0 = b$, such that $1 - S = (z^* - z)\,u$. Construct both
at once by induction on $y$-order. Order $0$ is $u_0(1/b - z) = 1 - bz$. At order
$j$, collecting $[y^j]$ of $(z^*-z)u = 1-S$ gives
$(1/b - z)\,u_j = N_j - b\,\zeta_j$ with
$N_j = -\sigma_j - \sum_{m<j}\zeta_m u_{j-m}$ known and of degree at most $j+1$;
choosing $\zeta_j := N_j(1/b)/b$ makes the right-hand side vanish at $z = 1/b$,
so the division is exact in $\Q[z]$ and $\deg u_j \le j$. Since $u(0,0) = b \ne 0$,
$u$ is a unit, and substituting $z = z^*$ shows $z^*$ is the unique root of
$1-S$ in $1/b + y\Q[[y]]$.

\paragraph{Step 2: the inverse unit stays polynomial per order.} Expanding
$u^{-1}$ as a geometric series in $\sum_j (u_j/b)y^j$, the $y^k$-coefficient
takes only $m \le k$ factors and compositions of $k$, so
$\deg [y^k]u^{-1} \le k$. Hence for $G := E_b E_t u^{-1}$, splitting
$k = k_b + k_t + k_u$ gives $\deg [y^k]G \le k+1$; equivalently, writing
$G = \sum_i g_i(y) z^i$, we have $\mathrm{ord}_y(g_i) \ge i-1$. This fact makes the
resummation exact rather than asymptotic.

\paragraph{Step 3: exactly one geometric mode.} Set $\mu := 1/z^*$. In
$\Q[[y]][[z]]$, $(z^*-z)^{-1} = \mu\sum_m (\mu z)^m$, so
\[
  [z^H](F - P) \;=\; \sum_{i \le H} g_i\, \mu^{\,H+1-i}
  \;=\; \mu^{\,H+1}\bigl(\hat C - \rho_H\bigr),
\]
with $\hat C := \sum_{i \ge 0} g_i (z^*)^i$ convergent $y$-adically and tail
$\rho_H$ of $y$-order at least $H$. For $H \ge k+1$ the tail contributes nothing
at order $y^k$, and $[z^H][y^k]P = 0$ as well since $\deg_z [y^k]P \le k < H$.
Setting $C := \hat C \mu$ gives the first display of Theorem~\ref{thm:C}. (Check:
$[y^0]G = z/b$, so $C(0) = 1/b$, and $k=0$ reads $T(H,H) = b^{\,H-1}$, the pure
walk.)

\paragraph{Step 4: production coordinates.} The second display reads the triangle
along fixed $n$ rather than fixed $H$, and the change of coordinates is the
diagonal form of Lagrange inversion. In the form needed here: for
$f, \varphi \in \Q[[w]]$ with $\varphi(0) \ne 0$ and $\hat w \in y\Q[[y]]$ the
unique solution of $\hat w = y\varphi(\hat w)$,
\[
  \sum_{k \ge 0} y^k\,[w^k]\bigl(f\varphi^k\bigr)
  \;=\; \frac{f(\hat w)}{1 - y\varphi'(\hat w)} ,
\]
proved by writing $[w^k](f\varphi^k) = \mathrm{res}_w(f\varphi^k w^{-k-1})$,
summing the geometric series to $\mathrm{res}_w\bigl(f/(w - y\varphi(w))\bigr)$,
and factoring $w - y\varphi(w) = (w-\hat w)v(w,y)$ by Step 1's induction with
$v(\hat w) = 1 - y\varphi'(\hat w)$. Applying it with $f = C\mu^n$ and
$\varphi = \mu^{-1}$ gives $\sum_k y^k [w^k](C\mu^{n-k}) = K(y)M(y)^n$, in which $n$
enters only through $M^n$; renormalising to the law's units turns that
into the exponential display, with
$a_j := [y^j]\log\bigl(b\,K(b^3 y)\bigr)$ and
$b_j := [y^j]\log\bigl(M(b^3 y)/b\bigr)$. Every object at order $j$ is built from
weights of surplus at most $j$, which is the locality claim. \qed

\begin{remark}[scope]
The proof above is written for a general row-local lattice, replacing the
constant $3$ of the king case by $b$ throughout. The king case is the one that
was written out first in full detail and the one formalized in Lean; the
generalization is the same induction with $b$ symbolic. Its two delicate steps,
the choice of $\zeta_j$ and the $y$-adic convergence of $\hat C$, use only
$b \ne 0$ and $\mathrm{ord}_y(g_i) \ge i-1$ respectively.
\end{remark}

\subsection{Two corollaries}

\begin{corollary}[two constants per level]
\label{cor:two}
$P_k(n) = \mathrm{known}_k(n) + a_k + b_k\,n$, where $\mathrm{known}_k$ is
determined by the levels below $k$. Consequently, given
Theorem~\ref{thm:A}'s degree bound, \emph{any two} in-onset values of a diagonal
determine $(a_k,b_k)$ by a nonsingular $2\times2$ rational solve, and every further
in-onset value is then a theorem rather than a test.
\end{corollary}

\begin{proof}
$[y^k]\exp(\sum_j c_j y^j)$ is a polynomial in $c_1,\dots,c_{k-1}$ plus $c_k$
itself; the new level enters linearly and $c_k = a_k + b_k n$.
\end{proof}

\begin{corollary}[leading coefficient]
\label{cor:lead}
The only contribution to $n^k$ at order $y^k$ is $(b_1 n y)^k/k!$, so
$[n^k]P_k = b_1^{\,k}/k!$. Since $P_1(n) = b_1 n + a_1$ and $[n]P_1 = \Wp$
(Definition~\ref{def:pair}), $[n^k]P_k = \Wp^{\,k}/k!$ and $\deg P_k = k$.
\end{corollary}

For the king lattice $P_1(n) = 25n - 45$ is derived from scratch by a
row-junction count, so $b_1 = 25$, $a_1 = -45$, and $[n^k]P_k = 25^k/k!$ for all
$k$.

\subsection{The defect gas}

Write
$F(n,u) = \sum_k P_k(n)u^k$ with $P_0 = 1$, and let $c_k = [u^k]\log F$. Then
Theorem~\ref{thm:C} says exactly that \emph{every $c_k$ is linear in $n$}: the
defects form a one-dimensional gas that is extensive, with all of the interaction
living in the constant terms. Degree $\deg P_k \le k$ follows because the top term
of $P_k$ is $c_1^k/k!$, and the leading coefficient is $\Wp^k/k!$ because $\Wp$
is the slope of $c_1$.

\begin{table}[H]
\centering
\begin{tabular}{l l l l l}
\toprule
lattice & $c_1$ & $c_2$ & $c_3$ & $c_4$ \\
\midrule
square & $4n-8$  & $-19n+54$              & $\tfrac{472}{3}n-\tfrac{1718}{3}$ & $-\tfrac{3099}{2}n+6558$ \\
hex    & $9n-15$ & $-\tfrac{37}{2}n-\tfrac{83}{2}$ & $32n-32$ & \\
king   & $25n-45$ & $-\tfrac{209}{2}n-\tfrac{891}{2}$ & & \\
\bottomrule
\end{tabular}
\caption{Cumulants of the diagonal series, computed by a drift-parametric
transfer computation (\repo{experiments/gas\_cumulants.py}). Linear in $n$ on
every lattice, at every level reached.}
\label{tab:cumulants}
\end{table}

The king $c_2$ is the calibration: the parametric computation returns the pair
$(a_2,b_2) = (-891/2,\,-209/2)$ that a king-only reconstruction supplies
independently.

\begin{table}[H]
\centering
\begin{tabular}{l r r r}
\toprule
cluster type & square & hex & king \\
\midrule
$(2)$       & 4  & 9   & 25 \\
$(3)$       & 9  & 16  & 49 \\
$(4)$       & 16 & 25  & 81 \\
$(2,2)$     & 12 & 60  & 339 \\
$(2,3)$     & 30 & 138 & 930 \\
$(2,2,2)$   & 36 & 409 & 4778 \\
\bottomrule
\end{tabular}
\caption{Interior cluster weights $W_c$, computed parametrically in the step set
$D$.}
\label{tab:weights}
\end{table}

The king column reproduces every entry of an independent king-only weight table.
Run at too small a search window the parametric version instead reproduces that
table's documented clipping bug, $(2,2,2) = 4776$, so the window size is checked
on every run.

\section{The universal spine}
\label{sec:spine}

The mod-$p$ behavior of the law is where quantifying over lattices pays off. The
answer is one curve; the lattice contributes only a prime and a scalar.

\begin{theorem}[universal spine]
\label{thm:B}
Let $p$ be a prime dividing $b$, and let $w := \Wp \bmod p$. Then modulo $p$ the
series $\mathcal H(u)$ of Theorem~\ref{thm:C} satisfies
\[
  \mathcal H \;=\; 1 + w\,u\,\mathcal H^{-2},
  \qquad\text{that is}\qquad
  \mathcal H^3 = \mathcal H^2 + w\,u \quad\text{over } \F_p .
\]
If $w \ne 0$ this is the cubic $W^3 = W^2 + t$ after the rescaling $t = wu$: the
curve does not depend on the lattice. If $w = 0$ the mod-$p$ triangle is trivial
in the band.
\end{theorem}

\begin{proof}
The chain identity gives a master equation
$\mathcal H = 1 + \sum_c \widehat W_c\, u^{k_c}\mathcal H^{-(k_c+\ell_c)}$ with
$\widehat W_c = W_c\, b^{\,2k_c - \ell_c - 1}$, after the substitution
$\mu = b\mathcal H$, $u = y\mu/b^3$. Now count valuations. For every cluster,
using $\ell_c \le k_c$ from \eqref{eq:elllek},
\[
  v_p(\widehat W_c) \;=\; (2k_c - \ell_c - 1)\,v_p(b) + v_p(W_c)
  \;\ge\; (k_c - 1)\,v_p(b) ,
\]
which is at least $1$ for every cluster except the one with
$k_c = \ell_c = 1$ --- the pair row. So modulo $p$ only the pair row survives,
with weight $\widehat W = \Wp$.
\end{proof}

\begin{remark}[the statement is exactly sharp]
Both branches are realised by ordinary lattices, so neither half of
Theorem~\ref{thm:B} is vacuous. $D = \{-2,0,2\}$ has $\Wp = 57 \equiv 0 \pmod 3$,
realising the degenerate branch; the interval set with $b = 5$ has
$\Wp = 114 \equiv 4 \pmod 5$, a unit but not $1$, so the scaling is genuinely
present and ``$w = 1$ always'' is false. Two further consequences of
Remark~\ref{rem:notsquares}'s closed form, for interval $D$: at odd $p \mid b$
both $b^3$ and $b(b+1)/2$ vanish, so $w \equiv 4$ --- always a unit, and the
spine for odd $p$ is $\mathcal H^3 = \mathcal H^2 + 4u$ with the $4$
lattice-independent too; at $p = 2$, $w \equiv \lfloor b/2 \rfloor \pmod 2$, so
the degenerate branch occurs exactly when $4 \mid b$.
\end{remark}

For the king lattice, $p = 3$ and $w = 25 \bmod 3 = 1$, so the spine is
$\mathcal H^3 = \mathcal H^2 + u$ over $\F_3$; for polyhexes, $p = 2$ and
$w = 9 \bmod 2 = 1$, giving the same cubic over $\F_2$. The square lattice has
$b = 1$ and no prime, so the theorem is vacuous there, correctly. What that one curve
is worth on a single lattice --- an automatic triangle, two boundary cases in
closed form, and the Smith normal form of the leading block --- is Part II.

\section{The law as a tool}
\label{sec:machine}

Theorem~\ref{thm:A} plus Corollary~\ref{cor:two} converts enumeration into proof.
Because the degree is at most $k$ and the onset is at $H \ge k+1$, it suffices to
enumerate $k+1$ in-onset values of a diagonal, or just two by
Corollary~\ref{cor:two}, to determine $q_k$.

A surplus-budgeted row transfer, parametric in the
step set $D$, carries as state the current row's occupied cells (normalized in
$x$), which component of the animal-so-far each belongs to, and the surplus spent
so far; components that lose their foothold in the current row are pruned as
unreconnectable. The cost therefore depends on $k$ and $H$, not on the number of
animals. Rows may have gaps, so positions are searched in a window, and every run
reports a stability check at window size and window size $+2$.

\begin{table}[H]
\centering\small
\begin{tabular}{l r r l l}
\toprule
lattice & $b$ & $k$ & $P_k(n)$ & status \\
\midrule
square & 1 & 1 & $4n-8$ & matches \oeis{A308359} (proved there) \\
square & 1 & 2 & $8n^2-51n+86$ & \textbf{A308359's conjecture} (Cor.~\ref{cor:a308359}) \\
square & 1 & 3 & $\tfrac{32}{3}n^3-140n^2+\tfrac{1960}{3}n-1090$ & no published form found \\
square & 1 & 4 & $\tfrac{32}{3}n^4-\tfrac{712}{3}n^3+\tfrac{12635}{6}n^2-\tfrac{52813}{6}n+14496$ & no published form found \\
hex & 2 & 1 & $9n-15$ & reproduces the computed $b=2$ form \\
hex & 2 & 2 & $\tfrac{81}{2}n^2-\tfrac{307}{2}n+71$ & no published form found \\
hex & 2 & 3 & $\tfrac{243}{2}n^3-774n^2+\tfrac{1897}{2}n+28$ & no published form found \\
king & 3 & 1 & $25n-45$ & reproduces the from-scratch derivation \\
king & 3 & 2 & $\tfrac{625}{2}n^2-\tfrac{2459}{2}n+567$ & reproduces the hand derivation \\
\bottomrule
\end{tabular}
\caption{Diagonal closed forms. Each is a theorem: Theorem~\ref{thm:A} supplies
the degree and the onset, and the tabulated values supply the coefficients.
Every row additionally survives 5--11 withheld values past the ones used to
determine it.}
\label{tab:machine}
\end{table}

Two rows of Table~\ref{tab:machine} are checks rather than results. The king
$k=2$ form was derived by hand, by a different argument, before this machinery
existed, and the parametric computation returns it coefficient for coefficient.
The square $k=3$ row can be checked against A308359's own printed triangle with
no computation at all: it gives $282$ at $n=7$ and $638$ at $n=8$, which are
that entry's values.

\section{A conjecture of A308359}
\label{sec:a308359}

OEIS \oeis{A308359}~\cite{a308359}, entered by R.~J.~Mathar in May 2019, is the
triangle of fixed polyominoes with $n$ cells by bounding-box \emph{width}, height
free --- which is the transpose of the square-lattice instance of $T(n,H)$. The
entry records $T(n,n-1) = 4n-8$ for $n \ge 3$ as known, and states:

\begin{quote}
\textbf{Conjecture:} $T(n,n-2) = 8n^2 - 51n + 86$ for $n \ge 5$.
\end{quote}

\begin{corollary}
\label{cor:a308359}
The conjecture holds.
\end{corollary}

\begin{proof}
This is Theorem~\ref{thm:A} at $b = 1$, $k = 2$. The theorem gives
$\deg P_2 \le 2$ and validity from $n \ge 2k+1 = 5$, exactly the range the
conjecture states, so three enumerated in-onset values determine the
polynomial and the conjecture is whatever those three values say. Enumerating
$n = 5, 6, 7$ gives $8n^2 - 51n + 86$.
\end{proof}

We ran that from our own enumeration rather than from the entry's numbers, in
three stages so that a disagreement would say where it lived: the row sums
against \oeis{A001168}; then the $k=1$ diagonal, which must reproduce the
$4n-8$ that Mathar records as \emph{proved}, an independent third party's check
that our triangle is the same object as theirs; and only then the quadratic. It reproduces every value
of the entry's printed triangle through $n = 11$.

Two controls, both of which must fail, and do
(\repo{experiments/oeis\_a308359\_check.py}):
\begin{itemize}
\item a \emph{linear} fit through $n = 5, 6$ must miss at $n = 7$, giving
  $105$ against the true $121$, or the degree is not under test;
\item the formula must \emph{fail} at $n = 2k = 4$, since the onset is sharp,
  giving $T(4,2) = 9$ where the quadratic says $10$, or the onset is not under
  test.
\end{itemize}

\section{Polyiamonds, where the theorem stops}
\label{sec:polyiamond}

Polyiamonds are the natural next lattice to try and they are \emph{not} in the
class. Cells are equilateral triangles; each has three edge-neighbors; and the
up/down orientation alternates along a row, so the adjacency is parity-dependent
rather than translation-invariant and hypothesis (U) fails --- there is no single
step set $D$ to count. This is not a technicality that a little more care would
remove. The counts are \oeis{A001420}: $2, 3, 6, 14, 36, 94, 250, 675, 1838,
\dots$, nothing like the $b=2$ instance's \oeis{A001207}
(Remark~\ref{rem:names}).

The conclusion nevertheless appears to survive, in the periodic form one would
guess. The indexing shifts, because a height-$H$ strip holds two triangle
orientations per column, so the \emph{maximal} $n$ at height $H$ is $2H-1$ rather
than $H$. Checked to $n = 12$ against a brute-force control
(\repo{experiments/universal\_law\_check.py}, \texttt{polyiamond\_probe}):
\[
  T(2H-2,\,H) = 2^{\,H-2},
  \qquad
  T(2H-1,\,H) = H\cdot 2^{\,H-1},
\]
and the next diagonal, $T(2H-2+k, H)/2^H$, is quadratic in $H$ at $k = 2$
(constant second differences, verified). The ground states at height $H$ are the
$2H-2$-cell up-down domino chains.

Making that a theorem means allowing a \emph{period-2} step set and asking for
$q_k$ quasi-polynomial with period $2$: rerun the five steps of
Section~\ref{sec:A} with a transfer over one period, so that the drift step
becomes a matrix. We have not done it, and we state it as an extension rather
than a result.

\section{Formalization}
\label{sec:formal}

The king instance is formalized in Lean~4 (toolchain v4.31.0) against
mathlib~\cite{mathlib}, in a development of 65 modules and about 21{,}700 lines
with \textbf{zero occurrences of \texttt{sorry}}. What is relevant here:

\begin{itemize}
\item The engine (\texttt{Universal.universal\_shape\_d},
  \texttt{universal\_shape}, \texttt{universal\_shape\_production},
  \texttt{universal\_production\_int\_all}) and the \texttt{kingSystem} and
  \texttt{peelSystem} instances depend on \textbf{standard axioms only}
  (\texttt{propext}, \texttt{Classical.choice}, \texttt{Quot.sound}).
\item The product form (\texttt{grand\_form} and \texttt{grand\_form\_prod},
  Theorem~\ref{thm:C}) likewise depends on standard axioms only. It was
  restructured for Lean as a \emph{staircase} route: in sequence form, Steps 1--3
  collapse to a convolution fixed point plus one strong induction, and Step 4's
  Lagrange lemma becomes a fixed point, with no power-series ring involved. Same
  theorem, different scaffolding.
\item \texttt{lead\_coeff\_25} (Corollary~\ref{cor:lead} for the king) adds a
  single named \texttt{native\_decide} leaf, which is the evaluation of one
  cluster weight.
\item The per-lattice instance determinations (\texttt{king\_P1\_*},
  \texttt{square\_P1\_*}, \texttt{hex\_P1\_*}) and the cross-family row-sum
  gates carry named leaves that amount to two triangle cells per lattice.
\item The $k \le 18$ production pinning
  (\texttt{P\textit{k}\_grand\_of\_banked}) is \emph{conditional on enumerated
  cells}: those cells are hypotheses of the Lean theorems, not proved in Lean.
\end{itemize}

Of the axiom footprints, 88 are \emph{guarded}: those modules wrap
\texttt{\#print axioms} in \texttt{\#guard\_msgs}, so a changed footprint is a
build error and building the project checks them. The remaining 148 print their
footprints but are informational. There is no \texttt{sorryAx} and no anonymous
\texttt{Lean.ofReduceBool}; every compiled-evaluation leaf is named, so the
trusted surface is enumerable rather than diffuse.

The development does \emph{not} certify the enumeration engine, the stored
triangle, or any anchor value. Those enter as explicit hypotheses or as named
leaves.

% ===========================================================================
\part*{Part II --- the mod-$3$ arithmetic of the king triangle}
\addcontentsline{toc}{part}{Part II --- the mod-$3$ arithmetic of the king triangle}

Part I proved that modulo a prime $p$ dividing $b$ the law has one spine,
$\mathcal H^3 = \mathcal H^2 + wu$, independent of the lattice
(Theorem~\ref{thm:B}). This part is what that curve is worth on a single
lattice. Taking the king instance, where $p = 3$ and $w = 25 \equiv 1$, the
whole triangle modulo $3$ turns out to be $3$-automatic, its boundary cases
close, and the Smith normal form of its leading block can be counted.

The headline is a statement about an integer matrix rather than about a
generating function. The leading $N \times N$ block $T_N$ is lower triangular,
since $T(n,H) = 0$ for $H > n$, and every invariant factor of its Smith normal
form over $\Z$ is a power of $3$. The number of \emph{nontrivial} ones, those
with exponent $> 0$, does not follow from the determinant:

\begin{table}[H]
\centering\small
\begin{tabular}{l r r r r r r r r r r r}
\toprule
$N$ & 4 & 5 & 6 & 7 & 8 & 9 & 10 & 11 & 12 & 13 & 14 \\
\midrule
nontrivial factors & 1 & 2 & 2 & 2 & 3 & 3 & 3 & 4 & 4 & 4 & 5 \\
\bottomrule
\end{tabular}
\caption{One new nontrivial invariant factor per three columns.}
\label{tab:snfcount}
\end{table}

\section{Setup}
\label{sec:setup}

Part I's Theorem~\ref{thm:A}, re-indexed by $n$ rather than $H$ and
specialized to the king lattice ($b = 3$, $\Wp = 25$), reads
\begin{equation}
  T(n,\,n-k) \;=\; P_k(n)\cdot 3^{\,n-1-3k}
  \qquad\text{for every } n \ge 2k+1,
  \label{eq:law}
\end{equation}
and Theorem~\ref{thm:C}, the product form, gives $G, H \in \Q[[y]]$ with
$G(0) = H(0) = 1$ and
\begin{equation}
  \sum_{k \ge 0} P_k(n)\,y^k \;=\; G(y)\,H(y)^n
  \qquad\text{for every } n \ge 2k+1 \text{ at order } y^k .
  \label{eq:grand}
\end{equation}
Equation \eqref{eq:grand} is what makes this part possible: it replaces
infinitely many polynomials by two power series. The cubic of
\S\ref{sec:cubic} is Theorem~\ref{thm:B} at $p = 3$, where the single-defect
weight is $25 \equiv 1 \pmod 3$, so no rescaling is visible.

\begin{proposition}[P\'olya integrality]
\label{prop:polya}
Each $P_k$ is integer-valued, hence $3$-integral: $\vthree(P_k(n)) \ge 0$ for
every integer $n$.
\end{proposition}

\begin{proof}
The window $[2k+1,\,3k+1]$ contains exactly $k+1$ consecutive integers, and on
it $P_k(n) = T(n,n-k)\cdot 3^{\,3k+1-n}$ is a nonnegative count times a
nonnegative power of $3$, hence an integer. A polynomial of degree $k$ that is
integral at $k+1$ consecutive integers is integer-valued
everywhere~\cite{polya1915}.
\end{proof}

The leading coefficient of $P_k$ is $25^k/k!$, so $P_k$ takes integer values
without having integer coefficients, as $\binom{n}{k}$ does; the window in the
proof is the $k+1$ consecutive integers such a polynomial needs.

\begin{proposition}
\label{prop:GH}
$G = \sum_k P_k(0)y^k$ and $H = \bigl(\sum_k P_k(1)y^k\bigr)/G$ are integer power
series with constant term $1$.
\end{proposition}

\begin{proof}
Evaluate \eqref{eq:grand} at $n = 0$ and $n = 1$. Both left sides are integral
by Proposition~\ref{prop:polya}; $G$ has constant term $P_0(0) = 1$, so it is
invertible in $\Z[[y]]$ by the geometric series.
\end{proof}

Proposition~\ref{prop:GH} computes $H$ from the diagonal polynomials, and with
$P_k$ known through $k = 17$ it begins $1, 25, 208, 1483, 20688, 130208, \dots$
It is the per-row transfer series of Part I's cluster decomposition, and its
first coefficient $25$ counts the single-defect species of one row, which is
Part I's pair weight $\Wp$. Whether $H$ is a known sequence is for the missing novelty search
to settle; it has not been looked up in the OEIS~\cite{oeis}.

\section{Three coefficient identities}
\label{sec:ladder}

Everything below uses three congruences, each derived from Part I's cluster master
equation
\[
  H \;=\; 1 + \sum_c \widehat W_c\, u^{k_c} H^{-(k_c+\ell_c)} ,
\]
the sum running over cluster types $c$ of surplus $k_c$ and length $\ell_c$.

\begin{description}
\item[$(\star a)$] $G \equiv 1 \pmod 9$, from a residue formula for the boundary
  clusters, a valuation lemma on boundary weights, and one identity of the
  mod-$9$ cubic.
\item[$(\star b)$] $H \equiv W \pmod 3$. Every cluster has surplus at least its
  length, so the valuation lemma leaves only the two-cell row surviving modulo
  $3$: the cubic counts chains of those rows and nothing else. The derivation
  holds at all orders.
\item[$(\star c)$] $(H^3 - H(t^3))/3 \equiv t^2 + tW \pmod 3$, from the mod-$9$
  cubic alone. Uniqueness of its solution gives
  $H(t^3) \equiv H^2 + 25t - 3t^2 - 3tW \pmod 9$, and a three-line cancellation
  over $\F_3$ finishes.
\end{description}

$(\star b)$ reproduces $342$ independently enumerated triangle cells through the
digit product of \S\ref{sec:cubic}, and $(\star c)$ holds to order $t^{300}$
against the mod-$27$ solution of \S\ref{sec:lift}.

\section{The spine cubic and the digit product}
\label{sec:cubic}

\begin{definition}
$W \in \F_3[[t]]$ is the unique root with $W(0) = 1$ of $W^3 = W^2 + t$.
\end{definition}
Written $W^2(W-1) = t$, the cubic determines the coefficients of
$W = 1 + \dots$ one after another, which gives existence and uniqueness; $W$ is
algebraic of degree $3$ over $\F_3(t)$. Every computation below goes through its
rational parametrisation
\begin{equation}
  t \;=\; \frac{x}{(1-x)^3},
  \qquad
  W \;=\; \frac{1}{1-x} ,
  \label{eq:param}
\end{equation}
checked by $W^3 - W^2 = (1-x)^{-3}\bigl(1 - (1-x)\bigr) = x(1-x)^{-3} = t$.

\begin{theorem}[digit product]
\label{thm:digit}
Write $n = \sum_i n_i 3^i$ in base $3$. Then
\[
  P_k(n) \bmod 3 \;=\; [y^k]\;\prod_i W\!\left(y^{3^i}\right)^{n_i}.
\]
\end{theorem}

\begin{proof}
By $(\star a)$ and $(\star b)$, the product form \eqref{eq:grand} reduces modulo
$3$ to $\sum_k P_k(n)y^k \equiv W^n$. Frobenius over $\F_3$ gives
$W(y)^{3^i} = W(y^{3^i})$, so $W^n = \prod_i \bigl(W^{3^i}\bigr)^{n_i} =
\prod_i W(y^{3^i})^{n_i}$.
\end{proof}

Since $W$ is algebraic over $\F_3(t)$, Christol's
theorem~\cite{christol1980} makes the array $P_k(n) \bmod 3$
$3$-automatic~\cite{allouche2003}, and Theorem~\ref{thm:digit} writes that
automaton out in two dimensions.

\section{The spine of a column, odd and even}

Call the cell $T\bigl(n, \lfloor 2n/3 \rfloor\bigr)$ at
$n = \lfloor 3H/2\rfloor$ the \emph{spine} of column $H$; it is the first entry
of that column not divisible by $3$. Odd and even $H$ need separate treatment,
because the exponent $n-1-3k$ of \eqref{eq:law} is $0$ at the spine in the first
case and $-1$ in the second.

Both proofs use Lagrange--B\"urmann in the following form~\cite[Ch.~5]{stanley2012}:
if $x = t\varphi(x)$ then
\begin{equation}
  [t^k]\,F(x(t)) \;=\; [x^k]\;F(x)\,\varphi(x)^k
  \left(1 - \frac{x\varphi'(x)}{\varphi(x)}\right),
  \label{eq:LB}
\end{equation}
valid over any commutative ring. With the parametrisation \eqref{eq:param},
$\varphi(x) = (1-x)^3$, so
\[
  \varphi'(x) \;=\; -3(1-x)^2 \;=\; 0 \quad\text{in characteristic } 3 .
\]
The correction factor in \eqref{eq:LB} is therefore identically $1$.

\begin{theorem}[odd spine]
\label{thm:oddspine}
$T(3k+1,\,2k+1) \equiv 1 \pmod 3$ for every $k \ge 0$.
\end{theorem}

\begin{proof}
At $n = 3k+1$ the exponent in \eqref{eq:law} is $n-1-3k = 0$, so
$T(3k+1,2k+1) = P_k(3k+1) \equiv [t^k]W^{3k+1} \pmod 3$ by
Theorem~\ref{thm:digit}. Apply \eqref{eq:LB} with $F = (1-x)^{-(3k+1)}$:
\[
  [t^k]W^{3k+1} \;=\; [x^k]\,(1-x)^{-(3k+1)}(1-x)^{3k}
  \;=\; [x^k](1-x)^{-1} \;=\; 1 .
  \qedhere
\]
\end{proof}

\begin{theorem}[even spine]
\label{thm:evenspine}
$P_k(3k) \equiv 3 \pmod 9$ for every $k \ge 1$; equivalently
$T(3k,2k) = P_k(3k)/3 \equiv 1 \pmod 3$.
\end{theorem}

\begin{proof}
Let $S := (H^3 - H(t^3))/3$, integral because $H^3 \equiv H(t^3) \pmod 3$ is the
classical Frobenius congruence for integer series.

Suppose first $3 \nmid k$. Then
$H^{3k} \equiv H(t^3)^k + 3k\,S\,H(t^3)^{k-1} \pmod 9$. The first term involves
only exponents divisible by $3$, so $[t^k]$ annihilates it. Using $(\star a)$,
$(\star c)$ and $W(t^3) = W^2 + t = W^3$,
\[
  P_k(3k) \;\equiv\; 3k\,[t^k]\bigl(t^2 + tW\bigr)W^{3k-3} \pmod 9 .
\]
By \eqref{eq:LB} again, whose correction term is $1$, together with
\eqref{eq:param},
\[
  [t^k]\bigl(t^2+tW\bigr)W^{3k-3}
  \;=\; [x^k]\left(\frac{x^2}{(1-x)^3} + \frac{x}{1-x}\right)
  \;=\; \binom{k}{2} + 1 ,
\]
and for $3 \nmid k$ one has
$k\bigl(\binom{k}{2}+1\bigr) \equiv k\cdot k \equiv 1 \pmod 3$, so
$P_k(3k) \equiv 3 \pmod 9$.

If $3 \mid k$, write $k = 3m$. The middle term now carries a factor
$3 \cdot 3m \equiv 0 \pmod 9$, leaving the self-similarity
\[
  P_{3m}(9m) \;\equiv\; [t^{3m}]H(t^3)^{3m}
  \;=\; [\tau^{m}]H(\tau)^{3m} \;\equiv\; P_m(3m) \pmod 9 ,
\]
and induction on $\vthree(k)$ reduces to the previous case.
\end{proof}

We have verified the self-similarity for $m \le 5$, and the congruence
$P_k(3k)\equiv3\pmod 9$ for all $k \le 17$, against the enumerated triangle.

\begin{theorem}[the activation boundary]
\label{thm:activation}
Column $H$ of the triangle satisfies $T(n,H) \equiv 0 \pmod 3$ for every
$n < \lfloor 3H/2\rfloor$, and $T\bigl(\lfloor 3H/2\rfloor, H\bigr)
\not\equiv 0 \pmod 3$.
\end{theorem}

\begin{proof}
With $k = n - H$, \eqref{eq:law} gives
$\vthree(T(n,H)) = (3H-2n-1) + \vthree(P_{n-H}(n))$.

\emph{Below the boundary.} $n < \lfloor 3H/2\rfloor$ makes the explicit exponent
$3H-2n-1 \ge 1$, and $\vthree(P_k(n)) \ge 0$ by Proposition~\ref{prop:polya}. So
$\vthree(T(n,H)) \ge 1$. This half needs only the law, not \S\ref{sec:ladder}.

\emph{At the boundary.} $H$ odd puts $n = 3k+1$, where the explicit exponent is
$0$ and $\vthree(P_k(3k+1)) = 0$ by Theorem~\ref{thm:oddspine}. $H$ even puts
$n = 3k$, where the explicit exponent is $-1$ and $\vthree(P_k(3k)) = 1$ by
Theorem~\ref{thm:evenspine}. Either way $\vthree(T) = 0$.
\end{proof}

\section{The Smith normal form}
\label{sec:snf3}

\begin{theorem}
\label{thm:snf-3power}
For every $N$, the Smith normal form of $T_N := [T(i,j)]_{i,j \le N}$ over $\Z$
is $\mathrm{diag}(3^{e_1},\dots,3^{e_N})$. The cokernel $\Z^N/T_N\Z^N$ is a
finite abelian $3$-group.
\end{theorem}

\begin{proof}
$T(n,H) = 0$ for $H > n$, so $T_N$ is lower triangular and
$\det T_N = \prod_n T(n,n) = \prod_n 3^{\,n-1} = 3^{N(N-1)/2}$, a pure power of
$3$. The invariant factors form a divisibility chain $d_1 \mid \dots \mid d_N$
with $\prod d_i = |\det|$, so any prime $q \ne 3$ dividing some $d_i$ would
divide $d_N$, hence the determinant --- impossible.
\end{proof}

The argument uses only the determinant, so Theorem~\ref{thm:snf-3power} is
unconditional. Counting the nontrivial factors needs Theorem~\ref{thm:activation}, and
with it the congruences of \S\ref{sec:ladder}.

\begin{theorem}
\label{thm:snf}
$T_N$ has exactly $\lceil (N-1)/3\rceil$ nontrivial invariant factors.
\end{theorem}

\begin{proof}
By Theorem~\ref{thm:snf-3power} the invariant factors are $3$-powers, so the
number of nontrivial ones is $N - \mathrm{rank}(T_N \bmod 3) =
\mathrm{nullity}(T_N \bmod 3)$.

By Theorem~\ref{thm:activation}, column $H$ of $T_N$ is identically zero
modulo $3$ if $\lfloor 3H/2\rfloor > N$, and otherwise has its first nonzero
entry in row $\lfloor 3H/2\rfloor$. The map $H \mapsto \lfloor 3H/2\rfloor$ is
strictly increasing, so the nonzero columns have their leading entries in
distinct, increasing rows: they are in echelon position and hence linearly
independent modulo $3$. The null space is therefore spanned exactly by the zero
columns, and
\[
  \mathrm{nullity}(T_N \bmod 3)
  \;=\; \#\{H \le N : \lfloor 3H/2\rfloor > N\}
  \;=\; \left\lceil \frac{N-1}{3} \right\rceil ,
\]
the last step elementary.
\end{proof}

We have verified this directly on the enumerated matrix for every $N \le 36$,
reproducing Table~\ref{tab:snfcount} and the full exponent lists.

\section{Reading along rows}

The results above are statements about columns. Read along rows, the same
triangle has a last nonzero entry in each row, which
Theorem~\ref{thm:activation} already locates for two residue classes of $n$
and not for the third.

\paragraph{The spine is also the last nonzero of its row.} Being to the right of
the spine in row $n$ means $n < \lfloor 3H/2\rfloor$, the same region
Theorem~\ref{thm:activation} proves zero, so the two statements are transposes
of one fact.

\paragraph{Rows $n \equiv 2 \pmod 3$ have no spine cell at all.} For those rows
$\lfloor 3H/2\rfloor$ never equals $n$, so the last nonzero entry is one step
earlier, at the deficit-$2$ cell, whose residue is $2$.

\begin{theorem}[deficit-$2$ residue]
\label{thm:deficit2}
$T(3m+2,\,2m+1) \equiv 2 \pmod 3$ for every $m \ge 0$.
\end{theorem}

\begin{proof}[Method]
Let $D(v) := \sum_m P_{m+1}(3m+2)\,v^m$ be the generating function of the
family, which is a diagonal of $G\,H^n$. Formal Lagrange--B\"urmann over
$\Z/27$, together with the parametrisation $v = u/H^3$, turns the claim
$D = 5 + 18v/(1-v)$ into a rational identity on the mod-$27$ master curve
$E(H,u) = 0$. Since $E$ is monic in $H$, that identity is decided by exact
polynomial division, and the remainder is identically zero.
\end{proof}

So the last-nonzero residue of row $n$ cycles $1, 1, 2$ with $n \bmod 3$: the two
spine theorems supply the first two cases and Theorem~\ref{thm:deficit2} the
third.

\section{Sleeve zeros}
\label{sec:sleeve}

A cell to the left of the spine in row $n$ carries a forced deficit
$d := 3k+1-n$, and $\vthree(P_k(n)) = d$ generically. A \emph{sleeve zero} is a
cell with $\vthree(P_k(n)) > d$. No formula for them is known, and what follows
is a census.

It is complete for the project's data, $n \le 40$. Valuations are read off the
enumerated triangle through $\vthree(P_k(n)) = \vthree(T(n,n-k)) + d$, so no
fitted $P_k$ is needed; fitting would have stopped the census at $k = 17$, while
the proved regime $n \ge 2k+1$ reaches further.

\begin{table}[H]
\centering\small
\begin{tabular}{r l r r}
\toprule
$n$ & sleeve zeros, as $k\!: d \to \vthree(P_k(n))$ & count & longest run \\
\midrule
19 & $7\!: 3\to4$, \; $9\!: 9\to10$ & 2 & 1 \\
37 & $13\!: 3\to5$, \; $16\!: 12\to13$, \; $17\!: 15\to18$ & 3 & 2 \\
38 & $16\!: 11\to13$ & 1 & 1 \\
39 & --- & 0 & 0 \\
40 & $16\!: 9\to10$, \; $18\!: 15\to18$, \; $19\!: 18\to21$ & 3 & 2 \\
\bottomrule
\end{tabular}
\caption{Sleeve zeros in the first row where two occur, and in the last four
rows of the data. The four record events (first row with two, first adjacent
pair at $n=30$, first three in total at $n=35$, first three in a row at $n=36$)
are final for this data: rows 37 and 40 tie the three-in-total record but nothing
reaches four.}
\label{tab:sleeve}
\end{table}

In the final rows the $k = 16$ diagonal spikes at $n = 36, 37$ and $38$, and
again at $n = 40$, skipping only $n = 39$, which is the one sleeve free of zeros
after $n = 31$. The excess $\vthree - d$ there reaches $3$, against a ceiling of
$5$ over all $n$ (at $k=8$, $n=17$).

\subsection{Deficit two and three}
\label{sec:deficit23}

The census above is a census because no unit formula was available past the
spine. Two more deficits now have one, by the same route: fix the tower of
congruences that $H$ and $G$ satisfy modulo a power of $3$, then divide.

At $d = 3$ the statement is
\begin{theorem}
\label{thm:d3}
For every $k \ge 3$,
\[
  P_k(3k-2) \;\equiv\; 27\,r_k \pmod{81},
  \qquad r_k = 2,\,0,\,1 \ \text{ for } k \equiv 0,\,1,\,2 \pmod 3 .
\]
Equivalently, on the deficit-$3$ family $(n,H) = (3k-2, 2k-2)$, where
$T(3k-2,2k-2) = P_k(3k-2)/27$: $T \equiv 2 \pmod 3$ for $k \equiv 0$,
$T \equiv 1 \pmod 3$ for $k \equiv 2$, and $\vthree(P_k(3k-2)) \ge 4$ for
$k \equiv 1$. Since $r_k$ is a unit for $k \not\equiv 1$,
\[
  \vthree\bigl(P_k(3k-2)\bigr) > 3 \iff k \equiv 1 \pmod 3 .
\]
\end{theorem}

The last line settles what the sleeve census could only tabulate: at deficit $3$
the sleeve zeros are exactly the $k \equiv 1 \pmod 3$ cells, for \emph{all} $k$,
rather than for the $k \le 11$ the data reached.

\begin{remark}[the frame this rests on]
Theorem~\ref{thm:d3} is proved on an explicit frame and the frame is not all
theorem. (A1) the diagonal-law shape theorem and the grand form
$P_k(n) = [u^k]G(u)H(u)^n$ --- both proved and Lean-checked. (A2) $H$ congruent
to the fixed point of the mod-$81$ master equation at all orders, derived from
the renewal formalism with completeness from a proved row bound, and
machine-verified against the computed $h_k$ for $k \le 17$. (A3) the corresponding
congruence for $G$, likewise derived and verified to $k \le 17$. (A4) eleven
cluster weights, finite integers from a stored table with independent low-order
enumeration cross-checks. Everything above those four is unconditional algebra:
an inversion lemma, one implicit differentiation, a units inventory, and an exact
polynomial division whose remainder vanishes identically mod $81$. So (A1) and
(A4) are solid, and the weight sits on (A2) and (A3), which are derivations
checked against data rather than proofs.
\end{remark}

Deficit $2$ closes the same way and was done first
(Theorem~\ref{thm:deficit2}); the finer valuations at
$k \equiv 1 \pmod 3$ ($7, 4, 4, 5$ at $k = 4, 7, 10, 13$) are invisible
modulo $81$ and belong to the higher levels of the tower.

\begin{openproblem}[the deficit-$d$ unit formulas, $d \ge 4$]
Theorems~\ref{thm:oddspine}, \ref{thm:evenspine} and~\ref{thm:d3} are the cases
$d = 0, 1, 3$ of a family, with $d = 2$ proved the same way: give the unit part
of $P_k(n)$ at forced deficit $d$ for $d \ge 4$. Each step costs one more level
of the tower ($d = 3$ needed mod $81$), and the open remainder of the
Witt-vector tower of \S\ref{sec:lift} is exactly this.
\end{openproblem}

\section{The 3-adic lift}
\label{sec:lift}

Two further $3$-adic levels of the cubic have been computed.

\begin{itemize}
\item $H^3 - H^2 \equiv 25t \pmod{27}$. The spine cubic lifts two further
  $3$-adic levels, the constant being exactly the pair weight $25$
  ($25 \equiv 7 \bmod 9$ and $\equiv 1 \bmod 3$, which is why the mod-$3$
  statement reads $W^3 = W^2 + t$ with no visible constant).
\item $\bigl(H^3 - H^2 - 25t\bigr)/27 \equiv t\bigl(W(t^3)-1\bigr) \pmod 3$. The
  level-$3$ correction is again $W$, so the same series reappears one level up.
\item Also found: $\bigl(H(t^3) - H^2 - t\bigr)/3 \equiv -(t + t^2 + tW)
  \pmod 3$.
\end{itemize}

These three are measured, not proved.

\section{The \texorpdfstring{$n = 40$}{n = 40} predictions}

The theorems above were stated before the enumeration reached $n = 40$, so the
last rows test them rather than fit them. All three predictions held, cell by
cell, against the final triangle:

\begin{itemize}
\item $T(37,25) \equiv 1 \pmod 3$ (odd spine, $k=12$), and every in-regime cell
  of row $37$ above the spine vanishes modulo $3$. \textbf{Confirmed} --- the
  spine cell is $\equiv 1$, and $T(37,H)$ for $H = 26,\dots,37$ together with
  $T(n,25)$ for $n < 37$ all vanish.
\item $T(39,26) \equiv 1 \pmod 3$ and $P_{13}(39) \equiv 3 \pmod 9$ (even
  spine). \textbf{Confirmed}.
\item Every remaining in-regime cell modulo $3$, via the digit product.
  \textbf{Confirmed}.
\end{itemize}

\section{Related machinery, and why it does not apply}
\label{sec:related}

A reader who works on congruences for combinatorial sequences will reach for two
tools here, and neither reaches this object.

Rowland and Yassawi~\cite{rowlandYassawi2015} compute, for a sequence whose
generating function is the diagonal of a rational power series, a finite
automaton giving the sequence modulo $p^\alpha$ --- which would settle statements
of the kind proved here mechanically. The hypothesis fails: companion paper L4
proves the height generating function is not D-finite, so it is not the diagonal
of a rational function, and no automaton of that provenance exists to compute.
The obstruction is a theorem of ours rather than an absence of effort.

The second tool is the family of methods for the mod-$3^k$ behavior of
sequences satisfying finite-depth recurrences with polynomial coefficients.
Those apply to a sequence given by such a recurrence; $T(n,H)$ is given by no
recurrence we know, which is the reason this paper works through the master
equation and its fixed point instead.

A literature-priority pass on the spine cubic, the digit product and the Smith
normal form count found no collision (2026-08-18,
\repo{docs/priority-passes-2026-08-18.md}).

% ===========================================================================
\section{Novelty}
\label{sec:novelty}

Part I's material was covered by the project's first six literature
searches and by the 2026-08-18 pass, which looked only at what postdated them;
neither found a collision on the law, the spine or the product form. The pass
over Part II is recorded in the same place
(\repo{docs/priority-passes-2026-08-18.md}) and is quoted here because its
searches were different.


\textbf{Searched, no collision.} The project's first six literature searches covered
the diagonal law, haruspicy and non-D-finiteness, the staircase squeeze,
HV-convex polyplets and the growth-constant upper bound, and none of them
touched this paper's objects. A seventh ran on 2026-08-18 and is recorded in
\repo{docs/priority-passes-2026-08-18.md}: congruences modulo powers of $3$ for
combinatorial sequences, the Krattenthaler--M\"uller method for mod-$3^k$
behavior of recursive sequences, Rowland--Yassawi automatic congruences for
diagonals of rational functions, Smith normal forms of combinatorial triangles,
and base-$3$ digit-product formulas. It found no collision on the spine cubic,
the digit product or the Smith normal form count. It did find one
adjacent method, Rowland--Yassawi, treated in \S\ref{sec:related}. The nearest
neighbors otherwise are Christol's theorem, which makes ``algebraic over
$\F_3$'' and ``$3$-automatic'' the same statement, and the theory of
Artin--Schreier curves, to which the cubic below belongs.

% ===========================================================================
\section{Open problems}
\label{sec:open}


\begin{openproblem}[closed forms for multi-row cluster weights]
Table~\ref{tab:weights} is enumerated. Is there a closed form for $W_c$ as a
function of the step set $D$ and the type $c$, generalising
Remark~\ref{rem:notsquares}'s $\Wp(b)$ for interval $D$?
\end{openproblem}

\begin{openproblem}[the cumulants from the weights, for general $b$]
Table~\ref{tab:cumulants} says every $c_k$ is linear in $n$ on every lattice, and
Theorem~\ref{thm:C} proves it. Assembling the two coefficients of $c_k$ for
$k \ge 2$ \emph{from the cluster weights} has been done only for the king
lattice. Since the target's shape is now known, what is missing is exactly the
combination of weights that produces the pair $(a_k, b_k)$.
\end{openproblem}

\begin{openproblem}[the periodic extension]
Prove Section~\ref{sec:polyiamond}'s periodic extension: for a lattice that is
row-local only after passing to a period-$q$ $x$-translation, $T$ obeys the
diagonal law with $b$ the per-period drift and $q_k$ quasi-polynomial of
period $q$.
\end{openproblem}

\begin{openproblem}[the individual exponents]
Theorem~\ref{thm:snf} counts the nontrivial invariant factors and the
determinant fixes their sum at $N(N-1)/2$, but the individual exponents
$e_i(N)$ are unknown. The mod-$3$ rank gives the count and says nothing about
the higher $3$-adic structure; the lifts of \S\ref{sec:lift} are the natural
attack.
\end{openproblem}

\begin{openproblem}[other primes]
Theorem~\ref{thm:B} gives one cubic per prime dividing the drift count,
which for polyhexes is $p = 2$, with the same curve. Does the polyhex triangle
have a mod-$2$ arithmetic as rich as this one, with a digit product, an
activation boundary and a Smith normal form count? Nothing here forbids it and
nobody
has looked.
\end{openproblem}

% ===========================================================================
\section{Reproduction}
\label{sec:repro}


Every computation quoted above is in this repository and is re-runnable from a
clean checkout.

\begin{table}[H]
\centering\small
\begin{tabular}{@{}l R{0.56\textwidth}@{}}
\toprule
claim & check \\
\midrule
chain identity, king & \repo{experiments/diagonal\_law\_proof\_check.py} \\
Theorem~\ref{thm:A}, square instance & \repo{experiments/universal\_law\_check.py} \\
Theorem~\ref{thm:A}, hex instance & \repo{experiments/hex\_gas.py} \\
Table~\ref{tab:machine} (all rows) & \repo{experiments/diagonal\_machine.py} \\
Table~\ref{tab:cumulants} & \repo{experiments/gas\_cumulants.py} \\
$\Wp$, two independent routes & \repo{experiments/universal\_pair\_weights.py} \\
Corollary~\ref{cor:a308359} and its two controls & \repo{experiments/oeis\_a308359\_check.py} \\
onset sharpness, $k \le 5$ ab initio & \repo{experiments/spine\_deeper.py} \\
onset sharpness, all $k$ (mod 3) & \repo{experiments/depth1\_sharpness.py} \\
Section~\ref{sec:polyiamond} & \repo{experiments/universal\_law\_check.py} \\
Section~\ref{sec:formal} & \texttt{cd polyplets \&\& lake build} \\
this paper's printed numbers & \repo{paper/verify\_l\_papers.py} \\
\bottomrule
\end{tabular}
\caption{Where each claim is checked.}
\label{tab:repro}
\end{table}

The label identifying which lattice a step set denotes is not taken on faith:
\repo{experiments/diagonal\_machine.py} brute-forces the animal counts to $n=8$
on every invocation and matches them against the OEIS sequences named in
Example~\ref{ex:instances}, so a mislabelled instance is a failure rather than a
silently wrong table.

Part II's own checks:


\begin{table}[H]
\centering\small
\begin{tabular}{@{}l R{0.56\textwidth}@{}}
\toprule
claim & check \\
\midrule
T1--T5 and the congruences & \repo{experiments/ternary\_spine.py} (15 checks) \\
the 342-cell digit-product mass check & \repo{experiments/ternary\_spine.py} \\
Theorem~\ref{thm:deficit2} & \repo{experiments/deficit2\_proof.py} \\
Table~\ref{tab:sleeve} & \repo{experiments/sleeve\_zero\_census.py} \\
the enumerated triangle & \repo{results/triangle.txt} \\
the underlying law and product form & Part I; \repo{docs/proofs/diagonal-law.md}, \repo{docs/proofs/grand-form.md} \\
the master equation and weights & \repo{results/defect-gas.md} \\
\bottomrule
\end{tabular}
\caption{Where each Part II claim is checked.}
\label{tab:repro2}
\end{table}

\bibliographystyle{plain}
\bibliography{shared/refs}

\end{document}
